
\chapter[Long-Time Linear]{high-order implicit-explicit
Runge-Kutta methods for linear and semilinear system}



\section{Introduction}\label{Sec:intro}
In this work, we investigate a numerical approach to solve the following problem. 
Let $V \subset H = H' \subset V'$ be a Gelfand triple of Hilbert spaces, 
where the superscript ' denotes the dual. 
Namely, the embedding $V\embed H$ is continuous and dense, and 
$$(u, v)_{V, V'} = (u,v)_H \quad \forall u \in H \embed V', v \in V \embed H,$$
where $(\cdot, \cdot)_{V, V'}$ is the duality pairing between $V'$ and $V$, 
and $(\cdot,\cdot)_H$ is the inner product on $H$. 

We consider an abstract parabolic initial value problem: 
find
$$u \in L^2((0,T);V) \subset H^1((0,T);V')\embed C([0,T];H)$$\
such that 
\begin{equation}\label{eq:abstract}
    \begin{cases}
    \partial_t u = \calA u + f(t) & 0 < t < T,\\
    u(0) = u_0  \in H
  \end{cases}
\end{equation}
where $\calA: V \to V'$ is a bounded linear operator (possibly non-selfadjoint) 
with the following property that:
$$\beta^{-1}\|u\|_V^2 \leq -(Du, u) \leq \beta \|u\|_V^2 \quad \forall u \in V, $$
$$ \lvert(Lu, v)\rvert \leq C \|u\|_V \|v\|_H, 
\forall u\in V, v \in H,$$
where $D = (\calA+\calA^*)/2, L = (\calA-\calA^*)/2$ are the symmetric and skew-symmetric 
part of the operator $\calA$. Furthermore, $D$ is negative definite. 

\subsection{Examples of the abstract problem}
As a multiphysics system, the Stokes-Darcy system is considered in this work which describes a moving fluid governed by the Stokes equations in a free-flow region $\Omega_S\subset\R^d$ and a flow in a neighboring  porous media region $\Omega_D\subset\R^d$. Darcy flow and Stokes flow interact through an interface denoted by $\Gamma$ as shown in Figure \ref{NS_Darcy_domain}. Applications of such a coupled system are ubiquitous in nature, including groundwater system \cite{HoppePortaVassilevski:2007,HouQiuHeGuoWeiBai:2016,LaytonTranTrenchea:2013}, petroleum extraction \cite{ChenSunWang:2014}, and so on.

% the coupling of flow in a porous media region $\Omega_D\subset\R^d$ and a free-flow region $\Omega_S\subset\R^d$ separated by an interface $\Gamma$ (as shown in Figure \ref{NS_Darcy_domain}), which has many applications in groundwater system \cite{JHou_XMHe_CGuo_MWei_BBai_1,RHoppe_PPorta_YVassilevski_1,WLayton_HTran_CTrenchea_1}, industrial filtrations \cite{NHanspal_AWaghode_VNassehi_RWakeman_1,VNassehi_1},  petroleum extraction \cite{TArbogast_MGomez_1,JChen_SSun_XWang_1}, and so on.
\begin{figure}[t]
\begin{center}
\begin{tikzpicture}
%\draw[-, thick] (0,0) -- (4,0) ;
\draw[-,] (0,0) -- (4,0) ;
\draw[-,] (4,0) -- (4,3) ;
\draw[-, ] (4,3) -- (0,3) ;
\draw[-,] (0,3) -- (0,0) ;
\draw[-] (0,1.5) -- (4,1.5) ;
\draw (2.1,0.7) node[left]{$\Omega_S$} ;
\draw (2.1,2.1) node[left]{$\Omega_D$} ;
\draw (3,1.3) node[left]{$\Gamma$} ;
%\draw[-Latex,] (0.5,1.5) -- (0.5,1.0) node[below]{$\,\,\,\,{\bm n}_D$} ;
\draw[-Latex,] (0.5,1.5) -- (0.5,2.0) node[above]{${\mathbf{n}}$} ;
\end{tikzpicture}
\end{center}
\caption{Computational domain of the Stokes--Darcy system}
%s $\Omega_D$ and $\Omega_S$, and the interface $\Gamma$}
\label{NS_Darcy_domain}
\end{figure}
This model consists of a parabolic equation
\begin{equation}\label{DS-1}
\begin{cases}
\partial_t \phi -\nabla\cdot(\kappa\nabla\phi)  = f_D
&\mbox{in}~~\Omega_D\times(0,T], \\
\phi = 0 &\mbox{on}~~\partial\Omega_D
\backslash\overline\Gamma\times(0,T] ,\\
-\kappa\nabla\phi\cdot \bfn = \u\cdot \bfn   & \mbox{on}~~\Gamma\times(0,T] ,\\
\phi(0)= \phi_0 & \mbox{in}~~\Omega_D,
\end{cases}
\end{equation}
which describes the Darcy flow in the porous media region $\Omega_D$ through the unknown hydraulic head $\phi$,
%(by Darcy's law, the fluid discharge rate is given by $-\kappa\nabla\phi$),
and an evolving Stokes equation
\begin{equation}\label{DS-2}
\begin{cases}
\partial_t \u - \nabla\cdot \T(\u, p )
= \mathbf{f}_S &\mbox{in}~~\Omega_S\times(0,T] , \\
\nabla\cdot\u=0 &\mbox{in}~~\Omega_S\times(0,T] , \\
\u  = 0 &\mbox{on}~~\partial\Omega_S
\backslash\overline\Gamma\times(0,T] ,\\
-\T(\u,p)\bfn = g \phi \bfn +\mu(\bfu-(\bfu\cdot\bfn)\bfn)
& \mbox{on}~~\Gamma\times(0,T] , \\
\u(0)=\u_0 & \mbox{in}~~\Omega_S,
\end{cases}
\end{equation}
which describes free flow in the region $\Omega_S$ through the fluid velocity $\bfu$, where $\T(\u, p)=2\nu \mathbb{D}(\u)-p \mathbb{I}$ denotes the stress tensor, in which $\mathbb{D}(\u)=\frac12(\nabla\u+(\nabla \u)^\top)$ is the deformation tensor and $\mathbb{I}$ is the $d\times d$ identity matrix. The physical parameters $\kappa$, $g$, $\mu$ and $\nu$ in this model are positive constants, and $f_D$ and $\mathbf{f}_S$ are given source terms.

Homogeneous Dirichlet boundary conditions will be imposed on 
outer boundaries separately, i.e., on 
$\partial\Omega_D\backslash\Gamma$ and 
$\partial\Omega_S\backslash\Gamma$. 
The interface conditions on $\Gamma$ in \eqref{DS-1} 
and \eqref{DS-2} represent conservation of mass and 
balance of force, respectively, where $\bfn$ denotes the 
unit normal vector on $\partial\Omega_S$ as shown in Figure 
\ref{NS_Darcy_domain}.

For the ease of error estimate, we will rewrite the Stokes-Darcy system as an equivalent abstract problem \cite{li2020long}.

Let $H= L^2(\Omega_D)\times L^2(\Omega_S)^d$ and $V=H^1_\Gamma(\Omega_D)\times {\bf H}^1_\Gamma(\Omega_S; {\rm div_0})$, where
\begin{align*}
&H^1_\Gamma(\Omega_D)=\{\varphi\in H^1(\Omega_D):\varphi=0\,\,\mbox{on}\,\,\partial\Omega_D\backslash\Gamma\},\\
&{\bf H}^1_\Gamma(\Omega_S; {\rm div_0})=\{{\v}\in H^1(\Omega_S)^d:
\nabla\cdot{\v}=0\,\,\mbox{in}\,\,\Omega_S\,\,\mbox{and}\,\,{\v}=0\,\,\mbox{on}\,\,\partial\Omega_S\backslash\Gamma\}.
\end{align*}
The weak formulation of \eqref{DS-1}-\eqref{DS-2} reads: find $(\phi,\u)\in L^2((0,T);V)\cap H^1((0,T);V') \hookrightarrow C([0,T];H) $
satisfying the following equations for all test functions $(\varphi,{\v})\in L^2((0,T);V)$:
%for all test functions $(\varphi,{\bm v},q)\in H^1(\Omega_D)\times H^1(\Omega_S)^3\times L^2(\Omega_S)$ :
\begin{align}
&(\partial_t\phi,\varphi)_{D}+(\kappa\nabla \phi,\nabla \varphi)_{D}
-(\bfu\cdot\bfn,\varphi)_{\Gamma}=(f_D,\varphi) \label{DS-weak-1} \\[5pt]
&(\partial_t\bfu,{\bfv})_{S}+(2\nu \D(\bfu),\D({\bfv}))_{S}
% - (p,\nabla\cdot{\bm v}) \nonumber \\
+(g\phi,{\bfv}\cdot\bfn)_{\Gamma}
+\mu(\bfu-(\bfu\cdot\bfn) \bfn,{\bfv}-({\bfv}\cdot\bfn) \bfn)_{\Gamma} \label{DS-weak-2} \\
&=(\mathbf{f}_S,{\bfv}) \nonumber ,
%&&\forall\, {\bm v} \in H^1(\Omega_S)^3, \\[5pt]
%&(\nabla\cdot\u,q)_S =0 &&\forall\,q \in L^2(\Omega_S) .
\end{align}
where $(\cdot,\cdot)_D$ is the pairing between $H^1_\Gamma(\Omega_D)'$ and $H^1_\Gamma(\Omega_D)$, $(\cdot,\cdot)_S$ is the pairing between ${\bf H}^1_\Gamma(\Omega_S; {\rm div_0})'$ and ${\bf H}^1_\Gamma(\Omega_S; {\rm div_0})$, and
$(\cdot,\cdot)_\Gamma$ is the inner product on $L^2(\Gamma)$.

Let the operators $A_1:H^1_\Gamma(\Omega_D)\rightarrow H^1_\Gamma(\Omega_D)'$, $A_2:{\bf H}^1_\Gamma(\Omega_S; {\rm div_0})\rightarrow {\bf H}^1_\Gamma(\Omega_S; {\rm div_0})'$, $B:{\bf H}^1_\Gamma(\Omega_S; {\rm div_0})\rightarrow H^1_\Gamma(\Omega_D)'$ and $B^*:H^1_\Gamma(\Omega_D)\rightarrow {\bf H}^1_\Gamma(\Omega_1; {\rm div_0})'$
%\begin{align*}
%&M:H^1(\Omega_D)\rightarrow H^1(\Omega_D)',\\
%&D:{\bf H}^1(\Omega_S; {\rm div_0})\rightarrow {\bf H}^1(\Omega_S; {\rm div_0})',\\
%&B:{\bf H}^1(\Omega_S; {\rm div_0})\rightarrow H^1(\Omega_D)',
%%&B_h:{\bf X}_h^{r+1}(\Omega_S)\rightarrow {\bf X}_h^{r+1}(\Omega_S),\\
%\end{align*}
be defined via duality by
\begin{align*}
&(A_1\phi,\varphi)_D=(\kappa\overline{\nabla\phi},\nabla \varphi)_{D}
\hspace{181pt} \forall\, \phi, \varphi\in H^1_\Gamma(\Omega_D), \\
&(A_2 \u ,{\v} )_S=(2\nu \overline{\D(\u )},\D({\v} ))_{S}\!+\!\mu(\overline{\u\! -\!(\u \cdot\n) \n},{\v}\!-\!({\v} \cdot\n) \n)_{\Gamma}
\hspace{10pt} \forall\,\u ,{\v} \in  {\bf H}^1_\Gamma(\Omega_S; {\rm div_0}) ,\\
&(B \u ,\varphi)_D=(\overline{\u \cdot\n},\varphi )_{\Gamma}
\hspace{33pt}\forall\,\u\in {\bf H}^1_\Gamma(\Omega_S; {\rm div_0}) \hookrightarrow L^2(\Gamma),\,\,\forall\, \varphi \in H^1_\Gamma(\Omega_D) \hookrightarrow L^2(\Gamma),\\
&(B^*\phi, {\v})_D=(\phi,\overline{{\v} \cdot\n} )_{\Gamma}
\hspace{33pt}\forall\,\phi \in H^1_\Gamma(\Omega_D) \hookrightarrow L^2(\Gamma),\,\,\forall\,{\v}\in {\bf H}^1_\Gamma(\Omega_S; {\rm div_0}) \hookrightarrow L^2(\Gamma) .
\end{align*}
Then the weak formulation \eqref{DS-weak-1}-\eqref{DS-weak-2} can be written as
\begin{align}
&\partial_t\phi  +A_1 \phi  - B\u = f_{D} , \label{DS-op1} \\
&\partial_t\u  +A_2 \u + g B^*\phi  = \f_{S} . \label{DS-op2}
\end{align}
By defining notations
\begin{align}
u=\left(\begin{array}{cc}
\phi \\
g^{-\frac12}\u \end{array}\right) ,
\quad
f=\left(\begin{array}{cc}
f_{D} \\
g^{-\frac12}\f_{S} \end{array}\right)
\quad\mbox{and}\quad
\calA
=-\left(\begin{array}{cc}
A_1 &  -g^\frac12B \\
g^{\frac12}B^* &   A_2\end{array}\right) ,
\end{align}
equations \eqref{DS-op1}-\eqref{DS-op2} are equivalently reformulated to the abstract parabolic initial value problem \eqref{eq:abstract}.


\subsection{Review on existing studies}

Many works have been done for the Stokes-Darcy system \eqref{DS-1}-\eqref{DS-2}.
Beavers and Joseph \cite{BeaversJoseph:1967} firstly conducted experimental works on this system and proposed the well-known Beavers-Joseph interfacial conditions and then many other mathematical investigations follow \cite{Ene:1975,JagerMikelic:1996,MikelicJager:2000,JagerMikelicNeuss:2001,LevySanchez:1975,MarusicMikelic:2000}.
Using the conditions, analysis of the coupled Stokes-Darcy system  can be found in \cite{BurmanHansbo:2007,DAngeloZunino:2009,DiscacciatiMiglio:2002,LaytonFriedhelmYotov:2003,Riviere:2005,RiviereYotov:2005} for the steady case and in \cite{CaoGunzburgerHuaWang:2010,Zunino:2002} for the time-dependent case.

Due to the various applications of Stokes-Darcy system, many different numerical methods are developed and analyzed, including domain decomposition methods \cite{DiscacciatiQuarteroni:2004,BoubendirTlupova:2013,CaoGunzburger:2011,ChenGunzburger:2011,DiscacciatiMiglio:2002,DiscacciatiQuarteroni:2007,FengHeWangZhang:2012,GunzburgerHeLi:2018,VassilevWangYotov:2014,QiuHeLiLin:2020}, Lagrange multiplier methods \cite{LaytonFriedhelmYotov:2003,BabuskaGatica:2010,GaticaMeddahiOyarzua:2008,HuangChenCai:2012}, discontinuous Galerkin methods \cite{KanschatRiviere:2010,LipnikovVassilevYotov:2014,Riviere:2005,RiviereYotov:2005}, multigrid methods \cite{ArbogastGomez:2009,MuXu:2008}, partitioned time-stepping methods \cite{KubackiMoraiti:2015,MuZhu:2010,ShanZheng:2013,ZhangRuiCao:2019}, coupled finite element methods \cite{CaoGunzburgerHu:2010,CamanoGatica:2015,KarperMardalWinther:2009,MarquezMeddahiSayas:2015} and many others \cite{ChenGunzburger:2016,ChenGunzburger:2013,GiraultVassilevYotov:2014,HeJiangQiu:2020}. In particular, 
Kubacki et al \cite{KubackiTran:2017} presented an overview of non-iterative partitioned methods for such a system. With a time-step restriction for stability, numerical schemes of both first-order and high-order partitioned methods were presented. Gunzburger et al in \cite{GunzburgerHeLi:2018} analyzed a parallel, non iterative, multiphysics domain decomposition method for decoupling the Stokes-Darcy model with multistep backward difference formula (BDF) for the time discretization. Optimal order $O(\tau^k)$ for the $k$-step BDF scheme were established in a general framework for any $k\le5$.
Chen et al in \cite{ChenGunzburger:2013} proposed two second-order-in-time implicit-explicit methods including $2$-step BDF and second-order Adams-Moulton-Bashforth method (AME$2$) in which coupling term in the interface conditions was treated explicitly and established for both schemes the unconditional and uniform-in-time stability. Error bound was derived with $O(\tau^2)$.  An improvement of this work was presented in \cite{ChenGunzburger:2016}, in  which a third-order in time AME algorithm was studied and uniform-in-time error estimate was derived.
%Note that all methods are implicit-explicit (IMEX) schemes. 
Recently, authors in \cite{LiWangZhou:2020} presented an implicit-explicit (IMEX) scheme with $k$-step BDF in time and finite element discretization in space. In that paper, the spatial differential operator $\calA$ was split into a symmetric part and an anti-symmetric part on which implicit and explicit schemes were applied respectively. A symmetrized and decoupled temporal $k$-step BDF scheme was presented and optimal long-time error bound $O(\tau^k+h^2)$ was derived.

We notice that aforementioned high-order-in-time works are conducted with multi-step methods and to our knowledge there is no temporal high-order single-step methods adopted for the coupled Stokes-Darcy system in the literature. This motivates us to apply IMEX Runge-Kutta method on the Stokes-Darcy equations aiming to achieve high-order convergence in time.

When a time-dependent partial differential equation (PDE) involves terms of different types, it is natural to employ different discretizations for them. IMEX time-discretization schemes are an example of such an idea. Runge-Kutta methods are a family of implicit-explicit iterative methods which were developed around in the 1900s by the German mathematicians Carl Runge and Wilhelm Kutta. In \cite{AscherRuuthSpiteri:1997}, authors developed Runge-Kutta-based IMEX schemes, among which diagonally implicit Runge-Kutta (DIRK) \cite{WannerHairer:1996} scheme is a popular one and we will adopt DIRK in our scheme.



\subsection{Our contributions and the organization of the paper}

In this work, the spatial differential operator is split into symmetric part and skew-symmetric part, on which implicit and explicit Runge-Kutta schemes are applied, respectively. Based on assumptions ({\bf P1})-({\bf P4}) on the scheme, we establish the long-time error bound for the proposed scheme. 

The rest of this paper is organized as follows.  In Section \ref{sec:IMEX}, we present implicit-explicit Runge-Kutta schemes with various stages and orders and an algorithm for searching coefficients of IMEX Runge-Kutta schemes. Theoretical results for error estimate are shown in Section \ref{sec:convergence}. Finally we present numerical examples in Section \ref{sec:examples} to verify the proposed theorems.



\section{Implicit-Explicit Runge-Kutta Method}\label{sec:IMEX}

\subsection{Time-Stepping Schemes}

To build a time discretization of the equation \eqref{eq:abstract},
we split the interval $(0, T)$ as
$$0 = t^0 < t^1 < \cdots < t^n = T,$$
with the mesh size $\tau := \max_{1\leq i \leq n}\tau_i$ for $\tau_i = t^i - t^{i-1}$.
Let $u^n$, $u^{n,i}$ are the approximations to $u(t^n)$ and $u(t^{n,i})$ respectively.

Following Duhamel's principle, the solution of equation \eqref{eq:abstract}
can be shown as

$$u(t^n) = e^{\tau_n D}u(t^{n-1}) + \int_0^{\tau_n}e^{(\tau_n-s)D}(Lu(t^{n-1}+s)+f(t^{n-1}+s))\rmd s$$

Then a framework of a single step scheme of approximating $u(t^n)$ reads:

\begin{align}\label{eq:SingleStep}
    u^n = \sigma(-\tau_n D)u^{n-1} + \tau_n \sum_{i=1}^mp_i(-\tau_n D)(L u^{n,i} + f^{n,i})
\end{align}
where $t^{n,i} = t^{n-1}  + c_i\tau_n$
with $c_i$ are distinct real numbers in $[0,1]$ for $i = 1, \dots, m$,
$u^{n,i}$ is generated from $u^{n-1}$, $f^{n,i} = f(t^{n,i})$. 
Here $\sigma(\lambda)$ and $\{p_i(\lambda)\}_{i=1}^m$ are rational functions.
For simplicity reasons, we assume that $\tau^n$ is of uniform mesh size, i.e. $\tau^i = \tau$ for $i = 1, \dots, n$.

A series of famous schemes are included in this analysis, including the Implicit-Explicit Runge--Kutta scheme
we are going to talk about.

\subsection{Implicit-Explicit Runge-Kutta schemes}

For the symmetric part $D$ in equation \eqref{eq:abstract},
we consider an $m$-stage diagonally implicit Runge--Kutta (DIRK) scheme
with coefficient $A = \{a_{ij}\}_{m\times m}$, $b = \{b_{i}\}_{i=1}^m$.
For the skew-symmetric part, we make use of an $m$-stage explicit scheme
with the same abscissae $\hat{c}=c = \{c_i\}_{i=1}^m$
and coefficient $\hatA = \{\hata_{ij}\}_{m\times m}$, $\hatb = \{\hatb_{i}\}_{i=1}^m$.
Thus, the IMEX-RK can be determined by the following Butcher tabular
% \begin{table}[h]
%     \centering
%     \begin{tabular}[h]{c|cccc|cccc}
%         $c_1$ & $a_{11}$ & $0$ & $\dots$ & $0$ & $\hata_{11}$ & $0$ & $\dots$ & $0$ \\
%         $c_2$ & $a_{21}$ & $a_{22}$  & $\dots$ & $0$ & $\hata_{21}$ & $\hata_{22}$ & $\dots$ & $0$ \\
%         $\vdots$ & $\vdots$ & $\vdots$  & $\ddots$ & $\vdots$ & $\vdots$ & $\vdots$ & $\ddots$ & $\vdots$ \\
%         $c_m$ & $a_{m1}$ & $a_{m2}$  & $\dots$ & $a_{mm}$ & $\hata_{m1}$ & $\hata_{m2}$ & $\dots$ & $\hata_{mm}$ \\
%         \hline
%            & $b_1$ & $b_2$ & $\dots$ & $b_m$ & $\hatb_1$ & $\hatb_2$ & $\dots$ & $\hatb_m$
%     \end{tabular}
%     \caption{Butcher tableau for Runge--Kutta scheme.}
%     \label{tab:Butcher}
% \end{table}



Then the IMEXRK scheme can be written as

\begin{equation} \label{eq:imex}
    \begin{cases}
        u^{n,0} = u^{n-1}\\
        u^{n,i} = u^{n-1} + \tau \sum\limits_{j=1}^ia_{ij}Du^{n,j} + \tau \sum\limits_{j=1}^i\hata_{ij}(Lu^{n,j-1}+f(t^{n, j-1})) \quad i = 1, \dots, m\\
        u^n = u^{n-1} + \tau \sum\limits_{j=1}^m b_jDu^{n,j} + \tau \sum\limits_{j=1}^m \hatb_j(Lu^{n,j-1}+f(t^{n, j-1})) \\
    \end{cases}
\end{equation}

Or, in vector form, let $U^n = \{u^{n,i}\}_{i=1}^m$,
$V^n =  \{v^{n,i}\}_{i=1}^m =  \{Lu^{n,i-1}+f(t^{n,i-1})\}_{i=1}^m$.
We can get
\begin{equation} \label{eq:ImexVec}
    \begin{cases}
        U^n = \bbone u^{n-1} + \tau A DU^n + \tau \hatA V^n,\\
        u^n = u^{n-1} + \tau b^\top D U + \tau b^\top V.
    \end{cases}
\end{equation}

% Ususlly, we can make $b_i$ and $\hatb_i$ be the last column of $A$ and $\hatA$,
% which means
% \begin{equation}\label{eq:LastRow}
%     b^\top = z^\top A, \hatb^\top =  z^\top \hatA.
% \end{equation}
% where $z = (0, \dots, 0, 1)^\top$.

To illustrate the stability, we will make some further assumptions
to the scheme.
\begin{itemize}
    \item[{\bf (P1)}] $\sigma(\lambda)$, $p_i(\lambda)$ are polynomials with respect to $\lambda$ for all $i=1, \dots, m$.
    $0 < \sigma(\lambda) \leq 1$.
    Besides, the numerator of $\sigma$ is a polynomial of which the degree is one order lower than its denominator,
    and the numerators of $p_i$ are polynomials of which the degree are lower than their denominators.
    \item[{\bf (P2)}] $A = \{a_{ij}\}_{m\times m}$ is lower-triangular, and its diagonal entries $a_{ii}$ are positive for $i\geq 1$.
    % \item[{\bf (P2)}] $A = \{a_{ij}\}_{m\times m}$ is lower-triangular.
    \item[{\bf (P3)}] $b_i$ and $\hatb_i$ be the last column of $A$ and $\hatA$, which means that
    \begin{equation}\label{eq:LastRow}
        b^\top = z^\top A, \hatb^\top =  z^\top \hatA.
    \end{equation}
    where $z = (0, \dots, 0, 1)^\top$.
    \item[{\bf (P4)}] With the definition
    $$A_\sigma = \begin{pmatrix}0 & 0 \\ 0 & A \end{pmatrix}, \hatA_\sigma = \begin{pmatrix}0 & 0 \\ \hatA & 0 \end{pmatrix}, b_\sigma =\begin{pmatrix}0 \\ b\end{pmatrix},\hatb_\sigma =\begin{pmatrix}\hatb\\0\end{pmatrix},
    c_\sigma = \begin{pmatrix}0 \\c\end{pmatrix},$$
    we have the following properties for given positive integer $k$
    $${b}_*^\top c^{l-1} = \frac{1}{l}, \text{\quad for $1\leq l \leq k$, ${b}_* \in\{b_\sigma, \hatb_\sigma\}$}$$
    $$z^\top \left(\prod_{i=1}^{l}A_{*i}\right) e = \frac{1}{l!}, \text{\quad for $1\leq l \leq k$, $A_{*i} \in\{A_\sigma, \hatA_\sigma\}$}$$
    \begin{equation*}
        z^\top\left(\prod_{i=1}^{r}A_{*i}\right) c^l = \frac{l!}{(l+r)!}, \quad \forall l\geq 0, r\geq 0, l+r\leq k, A_{*i} \in\{A_\sigma, \hatA_\sigma\}
    \end{equation*}
\end{itemize}

\begin{remark}
    In condition that assumption (P3) holds, the third relation of 
    assumption (P4) is a sufficient condition of the other relations. 
\end{remark}
Then equation \eqref{eq:ImexVec} gives us
$$(I - \tau AD)U^n = \bbone u^{n-1} + \tau \hatA V^n,$$
Take the inverse and employ the last relation of equation \ref{eq:ImexVec},

$$U^n = (I - \tau AD)^{-1}\bbone u^{n-1} + \tau (I - \tau AD)^{-1}\hatA V^n,$$
$$u^n = z^\top U^n = z^\top(I - \tau AD)^{-1}\bbone u^{n-1} + \tau z^\top(I - \tau AD)^{-1}\hatA V^n$$


Let $\sigma(s) = z^\top(I + sA)^{-1}\bbone$,
and $p_i(s)$ be the $i$-th entry of $z^\top(I + sA)^{-1}\hatA$
for $s > 0$.
Then the scheme can also be written as
\begin{equation}\label{eq:sigma}
    u^n = \sigma(-\tau D) u^{n-1} + \tau \sum_{i=1}^m p_i(-\tau D)(L u^{n,i-1} + f(t^{n, i-1})).
\end{equation}

Later on we will show some properties of $\sigma$ and $p_i$.


\begin{lemma}\label{thm:Polynomial}
    If a Runge--Kutta scheme \ref{eq:imex} can be rewritten into \ref{eq:sigma},
    and the assumption (P2) and (P3) are valid,
    then all these $\sigma(s)$ and $p_i(s)$ are rational functions with respect to $s$.
    Their shared denominator is of degree $m$,
    and their numerators are of degree at most $m-1$.
\end{lemma}

Proof:
To show this, the only point we should focus on is the cofactor matrix
of $I+sA$. It is well-known that
$$(I+sA)^{-1} = \det(I+sA)^{-1} (I+sA)^*,$$
where the star denotes the cofactor matrix.

Since $\sigma(s) = z^\top(I + sA)^{-1}\bbone$,
$p_i(s)$ are the $i$-th entry of $z^\top(I + sA)^{-1}\hatA$,
we can read directly that $\sigma$ and $p_i$
share the same denominator $\det(I+sA)$, and
its order is $m$.

By the definition of the cofactor matrix,
we can know that all the entries of the cofactor matrix are polynomials of order at most $m-1$,
So $p_i$and $\sigma$ have their numerators of degree at most $m-1$,
which completes the proof.

\begin{corollary}\label{thm:PiSigmaBound}
    If all the requirements of Lemma \ref{thm:Polynomial} are valid,
    and Assumption (P1) holds,
    then there exists positive constants $c_0$, $c_1$, $c_2$, s.t.
    $$ \sigma(\lambda) < 1/(1+c_0\lambda), \forall \lambda > 0, $$
    $$ \lambda \sigma(\lambda) < c_1, \forall \lambda > 0, $$
    $$ p_i(\lambda)\sigma(\lambda)^{-1} < c_2, \forall \lambda > 0. $$
\end{corollary}

We list some schemes on the Appendix which satisfy the above assumptions, 
together with how we can find the schemes. 





\section{Convergence Analysis for Linear IMEX-RK}\label{sec:convergence}

Before we illustrate the stability theorem, we will show a stability lemma for the stages, which may be used in the later theorem.
\begin{lemma}\label{thm:Stages}
    If a series of solutions satisfy
    $$\psi = \phi + c_0 \tau D\psi + \tau \sum\limits_{i=1}^{m}a_{i}Dv^i + \tau \sum\limits_{i=1}^{m}\hata_{i}Lv^{i},$$
    for some constant $c_0>0$, then there exists a constant $C$ such that

    $$\|\psi\|_H \leq C\Big( \|\phi\|_H + \sum_{i=1}^m \|v_i\|_H \Big),$$
    $$\|\psi\|_V \leq C\Big( \|\phi\|_V + \sum_{i=1}^m \|v_i\|_V \Big),$$
\end{lemma}
Proof: Take the inverse of $(I - c_0 \tau D)$, we can get

$$\psi = (I - c_0 \tau D)^{-1}\Big(\phi + \tau \sum\limits_{i=1}^{m}a_{i}Dv^i\Big) + \tau \sum\limits_{i=1}^{m}\hata_{i}(I - c_0 \tau D)^{-1}Lv^{i}.$$
Test with $\psi$ and we can get
$$\|\psi\|_H^2 \leq C\|\psi\|_H\Big(\|\phi\|_H + \sum_{i=1}^m \|v_i\|_H\Big) + \tau \sum\limits_{i=1}^{m} \hata_i\Big(Lv^{i}, (I - c_0 \tau D)^{-1} \psi\Big)$$
and
$$\tau \Big(Lv^{i}, (I - c_0 \tau D)^{-1} \psi\Big) \leq C \tau \|v^i\|_H \|(I - c_0 \tau D)^{-1} \psi\|_V \leq C \|v^i\|_H \|\psi\|_H.$$
which gives us the first relation.

To get the second relation, instead of test $\psi$, we will now test $-D\psi$ to the above equation. It turns to be
$$\|\psi\|_V^2 \leq C\|\psi\|_V\Big(\|\phi\|_V + \sum_{i=1}^m \|v_i\|_V\Big) - \tau \sum\limits_{i=1}^{m} \hata_i\Big(Lv^{i}, D(I - c_0 \tau D)^{-1} \psi\Big)$$
and
$$\tau \Big(Lv^{i}, D(I - c_0 \tau D)^{-1} \psi\Big) \leq  C\tau \|v^i\|_V \|D(I - c_0 \tau D)^{-1} \psi\|_H \leq C \|v^i\|_V \|\psi\|_V.$$
which gives us the second relation.

\begin{theorem}\label{thm:Stable}
    If $u^i$ and $u^{n,i}$ are the solutions generated by \ref{eq:imex} and \ref{eq:sigma},
    with $f=0$, and the assumption (P1) (P2) (P3) hold,
    then
    $$\|u^n\|_H^2  \leq \|u^{n-1}\|_H^2 , \forall n>1$$
    when $\tau \leq \tau^*$, 
    where the constant $\tau^*$ is only 
    related to the scheme and $\beta$ in equation \eqref{eq:beta}, 
    and not related to $u$ or $u^n$.
\end{theorem}
Proof: Test equation \ref{eq:sigma} with $\sigma^{-1}(-\tau D)u^n$, we can get

$$\big(u^n, \sigma^{-1}u^n\big) = \big(\sigma u^{n-1}, \sigma^{-1}u^n\big) +\tau \sum_{i=1}^m \big(p_iL u^{n,i-1}, \sigma^{-1}u^n\big)$$

Due to the symmetry of $D$, all the rational combinations are symmetric and commutable, so 

$$\big(u^n, \sigma^{-1}u^n\big) = \big( u^{n-1}, u^n\big) +\tau \sum_{i=1}^m  \big(L u^{n,i-1},p_i \sigma^{-1}u^n\big)$$

Since $L$ is skew-symmetric, $(\alpha Lw,w) = \alpha(Lw, w)=-\alpha(w, Lw) = 0$ for all $\alpha, w$,
so
$$\big(u^n, \sigma^{-1}u^n\big)  = \big( u^{n-1}, u^n\big) +
\tau \sum_{i=1}^m  \bigg(L\big(u^{n,i-1}-p_i(0)^{-1}p_i \sigma^{-1}u^n\big),p_i \sigma^{-1}u^n\bigg).$$

For our assumption, there exist an $c_0' = c_0 /\beta$ s.t. $0 < \sigma(s) <1/(1+c_0'\beta s)$, so the LHS
of the above equation can be bounded by

$$\big(u^n, \sigma^{-1}u^n\big) > \big\|u^n\big\|_H^2 + c_0'\tau\big\|u^n\big\|_V^2$$

and
$$\big(u^{n-1}, u^n\big) = \frac12\big\|u^n\big\|_H^2+\frac12\big\|u^{n-1}\big\|_H^2-\frac12\big\|u^n-u^{n-1}\big\|_H^2.$$

For the second term on the RHS, use the property of $L$, we can get

\begin{align*}
\sum_{i=1}^m & \;\bigg(L\big(u^{n,i-1}-p_i(0)^{-1}p_i \sigma^{-1}u^n\big),p_i \sigma^{-1}u^n\bigg)\\
&\leq \epsilon \sum_{i=1}^m\big\|p_i \sigma^{-1}u^n\big\|_V^2
+ C_\epsilon \sum_{i=1}^m\big\|u^{n,i-1}-p_i(0)^{-1}p_i \sigma^{-1}u^n\big\|_H^2\\
&\leq \epsilon \sum_{i=1}^m\big\|p_i \sigma^{-1}u^n\big\|_V^2
+ 2C_\epsilon \sum_{i=1}^m\big\|u^{n,i-1}-u^n\big\|_H^2 +2C_\epsilon \sum_{i=1}^m\big\|u^n-p_i(0)^{-1}p_i \sigma^{-1}u^n\big\|_H^2\\
&=I_1+I_2+I_3
\end{align*}

By corollary \ref{thm:PiSigmaBound}, we can get
$$I_1 = \epsilon \sum_{i=1}^m\big\|p_i \sigma^{-1}u^n\big\|_V^2 < \epsilon_1 C \big\|u^n\big\|_V^2.$$

% The estimation of $I_2$ is more tricky.
% Reformulate the first relation of \ref{eq:ImexVec}
% and we can get
% \begin{align}\label{eq:ProofStable2}
%     U^n - u^n\bbone  = (u^{n-1}-u^n)\bbone  + \tau A D(U^n-u^n\bbone) + \tau \hatA (V^n-Lu^n\bbone )
% + \tau(AD+\hatA L)u^n \bbone.
% \end{align}
% Test it with $U^n - u^n\bbone$.
% $$\big(U^n - u^n\bbone\big)^\top \big(u^{n-1}-u^n\big)\bbone \leq \epsilon_{21} \sum_{i=1}^m \big\|u^{n,i}-u^n\big\|_H^2 + C_{21}\big\|u^{n-1}-u^n\big\|_H^2$$
% $$\big(U^n - u^n\bbone\big)^\top\tau A D\big(U^n-u^n\bbone\big) \leq \tau C_A \sum \big\|u^{n,i}-u^n\big\|_V^2$$
% $$\big(U^n - u^n\bbone\big)^\top\tau \hatA \big(V^n-Lu^n\bbone \big)\leq \tau C_\hatA \sum \big\|u^{n,i}-u^n\big\|_V^2 + \tau C_\hatA \sum \big\|u^{n,i}-u^n\big\|_H^2$$
% $$\big(U^n - u^n\bbone\big)^\top\tau AD u^n \bbone \leq \epsilon_{22} \big\|u^n\big\|_V^2 +  \tau^2 C_{22}\sum \big\|u^{n,i}-u^n\big\|_V^2 $$
% $$\big(U^n - u^n\bbone\big)^\top\tau \hatA L u^n \bbone \leq \epsilon_{23} \big\|u^n\big\|_V^2 +  \tau^2 C_{23}\sum \big\|u^{n,i}-u^n\big\|_H^2 $$

% Combine all these terms, for fixed small $\epsilon_{2i}$, and make $\tau < \tau_1^*$
% which is only related to $\epsilon_{2i}$ and $\hatA$,
% then the terms $\|u^{ni}-u^n\|^2$ of the RHS can partially cancel the LHS, then
% $$I_2 \leq C_1\epsilon_2 \big\|u^n\big\|_V^2 + C_1C_2\big\|u^{n-1}-u^n\big\|_H^2
% +\tau CC_1 \sum\big\|u^{n,i}-u^n\big\|_V^2$$

Then we are going to estimate $I_2$. The second relation in \ref{eq:imex} gives us
$$u^{n,i}-u^n = (u^{n-1}-u^n) + \tau \sum\limits_{j=1}^ia_{ij}D(u^{n,j}-u^n) + \tau \sum\limits_{j=1}^i\hata_{ij}L(u^{n,j-1}-u^n) + \tau c_i(D+L)u^n.$$

Define $w^{n,i} = u^{n,i}-u^n$, so
$$w^{n,i} = w^{n,0} + \tau \sum\limits_{j=1}^ia_{ij}Dw^{n,j} + \tau \sum\limits_{j=1}^i\hata_{ij}Lw^{n,j-1} + \tau c_i(D+L)u^n.$$

Similar with Lemma \ref{thm:Stages}, taking the inverse of $\big(I-a_{ii}\tau D\big)$, we can derive that

\begin{align*}w^{n,i} &= \tau \sum\limits_{j=1}^{i-1}a_{ij}\big(I-a_{ii}\tau D\big)^{-1}Dw^{n,j} 
+ \tau \sum\limits_{j=1}^i\hata_{ij}\big(I-a_{ii}\tau D\big)^{-1}Lw^{n,j-1} \\
&+ \big(I-a_{ii}\tau D\big)^{-1}w^{n,0} 
+ \tau c_i\big(I-a_{ii}\tau D\big)^{-1}Du^n+ \tau c_i\big(I-a_{ii}\tau D\big)^{-1}Lu^n. 
\end{align*}



Test with $w^{n,i}$, we can derive that 
\begin{align*}\|w^{n,i}\|_H^2 &\leq C \sum\limits_{j=1}^{i-1}a_{ij}\|w^{n,j}\|_H^2 
+ \tau \sum\limits_{j=1}^i\hata_{ij}\big(Lw^{n,j-1},\big(I-a_{ii}\tau D\big)^{-1} w^{n,i} \big)\\
&+ \|w^{n,0}\|_H\|w^{n,i}\|_H + C\tau^{1/2} \|u^n\|_V\|w^{n,i}\|_H+ \tau c_i\big(Lu^n, \big(I-a_{ii}\tau D\big)^{-1}w^{n,i}\big),
\end{align*}
and
$$\big(L\phi,\big(I-a_{ii}\tau D\big)^{-1} \chi \big) 
\leq C\|\phi\|_H\|\big(I-a_{ii}\tau D\big)^{-1} \chi\|_V
\leq C\|\phi\|_H\|\chi\|_H,$$
for any $\phi, \chi$. 
Therefore
$$\|w^{n,i}\|_H \leq C \|w^{n,0}\|_H + C \sum\limits_{j=1}^{i-1}a_{ij}\|w^{n,j}\|_H+C \tau^{1/2}\|u^n\|_V.$$

Accumulate from $1$ to $i$, we can derive that
$$\|w^{n,i}\|_H \leq C \|w^{n,0}\|_H + C \tau^{1/2}\|u^n\|_V,$$
which means that
$$\|u^{n,i}-u^n\|_H^2 \leq C \|u^{n-1}-u^n\|_H^2 + C \tau\|u^n\|_V^2.$$

For $I_3$, since $p_i\sigma^{-1}$ is bounded, we know that $\Big(1 -\frac{p_i(s)}{p_i(0)\sigma(s)}\Big)^2$
is also bounded for $s\geq 0$. Further, it is a rational function with value $0$ when $s = 0$, so
$$\bigg(1 -\frac{p_i(s)}{p_i(0)\sigma(s)}\bigg)^2 \leq Cs, \forall s \geq 0.$$
which means that
$$I_3 \leq C_\epsilon\tau \big\|u^n\big\|_V^2.$$

Now, we can get that

\begin{align*}
    \big\|u^n\big\|_H^2 =& - c_0\tau\big\|u^n\big\|_V^2 + \frac12\big\|u^n\big\|_H+\frac12\big\|u^{n-1}\big\|_H^2-\frac12\big\|u^n-u^{n-1}\big\|_H^2\\
    &+ \tau \epsilon C \sum_{i=1}^m\big\|u^n\big\|_V^2 + C_\epsilon\tau \|u^{n-1}-u^n\|_H^2 + C_\epsilon \tau^2\|u^n\|_V^2 \\
    &+ C_\epsilon\tau^2 \big\|u^n\big\|_V^2
\end{align*}

For a given IMEX--RK scheme, $\epsilon$ can be fixed so that $C_\epsilon$ is also a fixed value,
which is non-related to the time step $\tau$ and solution $u$.
Then a small $\tau$ can guarantee that
$C_\epsilon\tau \|u^{n-1}-u^n\|_H^2$ and $C_\epsilon\tau^2 \big\|u^n\big\|_V^2$ are bounded by the negative terms.
So finally we find

$$\big\|u^n\big\|_H^2 \leq  \big\|u^{n-1}\big\|_H^2, $$
which agrees with our claimant.\qed

% Lastly,
% \begin{align*}
%     \big\|u^{n,i}-u^n\big\|_V^2 &\leq 2\big\|u^{n,i}\big\|_V^2 + 2\big\|u^{n}\big\|_V^2\\
%         & \leq C\big\|u^{n-1}\big\|_V^2 + 2\big\|u^{n}\big\|_V^2
% \end{align*}
% where we use the result in Lemma \ref{thm:Stages} recursively. Then when $\tau \leq \tau_3^*$, we can get
% $$\big\|u^n\big\|_H^2 + \gamma \tau \big\|u^n\big\|_V^2 \leq \big\|u^{n-1}\big\|_H^2 + \gamma \tau  \big\|u^{n-1}\big\|_V^2.$$

% Let $\tau^* = \min \{\tau_1^*, \tau_2^*, \tau_3^*\}$, then
% we can have the above analysis. Note that
% $\tau^*$ now is only related to the scheme itself,
% and non-related to the solution $u$. A smaller $c_1$ requires a smaller $\tau^*$.


% \begin{remark}
%     If the matrix $A = \{a_{ij}\}$ is positive definite in sense of inner product,
%     i.e. $(A+A^T) \succ 0$,
%     then the second term of $I_2$ is negative definite,
%     and can cancel the high order terms when $\tau$ is small.
%     After that, the final stability will be $\big\|u^n\big\|^2 \leq \big\|u^{n-1}\big\|^2$
%     and the high regulared term will disappear.
% \end{remark}

Before the error estimation, 
we will show the convergence theorem for linear symmetric problem. 

\begin{theorem}
    A stiffly accurate scheme \eqref{eq:imex} 
    is accurate of order $k$ for linear symmetric problem
    if and only if 
\begin{subequations}
\begin{equation}
    z^\top \As^l \bbone = \frac{1}{l!}, \forall 0 \leq l \leq k. 
\end{equation}
and
\begin{equation}
    z^\top\As^r\Ahs c^l = \frac{l!}{(l+r+1)!}, \quad \forall l\geq 0, r\geq 0, l+r\leq k-1.
\end{equation}
\end{subequations}
\end{theorem}
Proof: let $U^n = \{u^{n,i}\}_{i=0}^m$,
$F^n =  \{f(t^{n,i})\}_{i=0}^m$.
We can get
\begin{equation} \label{eq:AppendImexVec}
    \begin{cases}
        U^n = \bbone u^{n-1} + \tau \As DU^n + \tau \Ahs F^n,\\
        u^n = u^{n-1} + \tau \bs^\top D U + \tau \bs^\top F.
    \end{cases}
\end{equation}

Here we write the tabular in extended form. 
Define a vector $z$ so that 
$$z^\top = (0, 0, \dots, 0, 1).$$
We know that a Runge--Kutta scheme is \emph{stiffly accurate} if and only if 
$$\bs^\top = z^\top \As, \bhs^\top = z^\top \Ahs.$$

The first relation of \eqref{eq:AppendImexVec} gives us that 
$$(I-\tau \As D)U^n = \bbone u^{n-1} + \tau \Ahs F^n, $$
$$U^n = (I-\tau \As D)^{-1}\bbone u^{n-1} + \tau (I-\tau \As D)^{-1}\Ahs F^n.$$
Take the last row of the above equation, which is the stiffly accuracy, 
so finally we gan get
$$u^n = \sigma(-\tau D)u^{n-1} + \tau \sum_{i=0}^mp_i(-\tau D)f(t^{n,i}). $$
where $\sigma(\lambda) = z^\top (1+\lambda \As)^{-1}\bbone$, 
$\{p_i(\lambda)\}^\top = z^\top (1+\lambda \As)^{-1} \Ahs $
for $\lambda \geq 0$. 
we know the rational real function on the operator $D$ is 
well-defined since $D$ is negative definite such that $-\tau D$ is positive definite. 

We know that the scheme \eqref{eq:imex} is accurate of order $k$ 
for linear symmetric problem if and only if 


\begin{subequations}
\begin{equation}\label{eq:Append_Acc1}
    \sigma(\lambda) = e^{-\lambda} + O(\lambda^{k+1}), \quad as~\lambda \to 0, 
\end{equation}
and, for any $0\leq l \leq k$, 
\begin{equation}\label{eq:Append_Acc2}
    \sum_{i=0}^m c_i^lp_i(\lambda) = \frac{l!}{(-\lambda)^{l+1}}
    \Big(e^{-\lambda} -\sum_{l=0}^j \frac{(-\lambda)^l}{l!}\Big)
    + O(\lambda^{k-l}), \quad as~\lambda \to 0. 
\end{equation}
\end{subequations}
See \cite{thomee2007galerkin}.

To get the first relation, since the rational function and 
exponential function are sufficiently smooth, we can just take the derivative and
test whether they are same at zero. 
$$\left(\frac{\rmd}{\rmd t}\right)^l\sigma = z^\top \left(\left(\frac{\rmd}{\rmd t}\right)^l (1+\lambda \As)^{-1}\right)\bbone
=(-1)^l l! z^\top (1+\lambda \As)^{-(l+1)}\As^l\bbone. $$
Therefore, \eqref{eq:Append_Acc1} is true if and only of
$$\left(\frac{\rmd}{\rmd t}\right)^l\sigma(0) = (-1)^l l! z^\top \As^l \bbone = (-1)^l, \quad \forall 0\leq l \leq k.$$

To evaluate the second relation, 
first we need some simplification.
The equation \eqref{eq:Append_Acc1} is equivalent to 
\begin{equation}\label{eq:Append_Acc2b}
    (-\lambda)^{l+1}\sum_{i=0}^m c_i^lp_i(\lambda) = l!
    \Big(e^{-\lambda} -\sum_{l=0}^j \frac{(-\lambda)^l}{l!}\Big)
    + O(\lambda^{k+1}), \quad as~\lambda \to 0. 
\end{equation}
Similarly, we only need to compare their derivatives at zero. 
Let $LHS =  (-\lambda)^{l+1}\sum_{i=0}^m c_i^lp_i(\lambda), RHS = l!
\Big(e^{-\lambda} -\sum_{l=0}^j \frac{(-\lambda)^l}{l!}\Big)$.

It is obviously that $LHS$ has a $\lambda^{l+1}$ factor 
so it is zero for no more than $l$ times derivative. 
So does $RHS$ because the second term is 
the Taylor's expansion of the first term. 
So we only need to test their derivatives with order 
more than $l$.

Let $r\geq 0$ and $l+1+r \leq k$. Take the $(l+1+r)$-th order 
derivative on $LHS$, with the Leibniz product rule, 
we can derive that
\begin{align*}
{\left(\frac{\rmd}{\rmd t}\right)}^{l+1+r}LHS(0)
&= C_{l+1+r}^r {\left(\frac{\rmd}{\rmd t}\right)}^{l+1} (-\lambda)^{l+1} {\left(\frac{\rmd}{\rmd t}\right)}^r\sum_{i=0}^m c_i^lp_i(\lambda) \\
&= \frac{(l+1+r)!}{(l+1)!r!}(-1)^{l+1}(l+1)!\cdot
z^\top (-1)^r r! (1+ 0 \cdot  \As)^{-(r+1)} \As^r \Ahs c^l\\
&= (-1)^{l+1+r}(l+1+r)!z^\top\As^r\Ahs c^l.
\end{align*}
    
Therefore, \eqref{eq:Append_Acc2} is true if and only of
$$(-1)^{l+1+r}(l+1+r)!z^\top\As^r\Ahs c^l = (-1)^{l+1+r}(l)!, \quad \forall l\geq 0, r\geq 0, l+r\leq k-1. $$
\qed

\begin{remark}
    If the source term is partially or fully computed implicitly, 
    then an additional requirement should be added that 
    \begin{equation*}
        z^\top\As^r\As c^l = \frac{l!}{(l+r+1)!}, \quad \forall l\geq 0, r\geq 0, l+r\leq k-1.
    \end{equation*}
The proof is a line-by-line copy of the previous. 
\end{remark}

Next, we shall derive an error estimate for the scheme. 
\begin{theorem}
    Suppose that the Assumption (P1), (P2), (P3) and (P4) are valid,  
    $u$ is the solution of \ref{eq:abstract}, $u^n$ is the solution of \ref{eq:imex}. 
    Then we can have the following error estimate

    $$\|u(t^n) - u^n\|_H \leq C\tau^k$$
    when $\tau \leq \tau^*$, 
    where the constant $\tau^*$ is only 
    related to the scheme and $\beta$ in equation \eqref{eq:beta}, 
    and not related to $u$ or $u^n$.
\end{theorem}
Proof: To begin with, we shall examine the truncation error for our scheme. 
Let $U_*^n = \{u(t^{n,i}\}_{i=0}^{m}$ be the exact solution. 
We define the local truncation error $R^n$ as 


\begin{equation}\label{eq:LocTrun}
U_*^n = \bbone u(t^n) + \tau \As D U_*^n + \tau \Ahs (L U_*^n + F^n) + R^n. 
\end{equation}
together with 
\begin{equation}
U^n = \bbone u^n + \tau \As D U^n + \tau \Ahs (L U^n + F^n). 
\end{equation}

Define $e^{n,i} = u(t^{n,i}) - u^{n,i}$, and vector $E^n$ be consisted with $e^{n,i}$. we can get

$$E^n = \bbone e^n + \tau \As D E^n + \tau \Ahs L E^n + R^n. $$

Divide $E^n$ into two vectors, such that 

$$E_1^n = \bbone e^n + \tau \As D E_1^n + \tau\Ahs L E_1^n, $$
$$E_2^n =  \tau \As D E_2^n + \tau\Ahs L E_2^n + R^n.$$
Obviously $E^n = E_1^n + E_2^n$, and 
$\|z^\top E_1^n \|_H \leq \|e^n\|_H$ by theorem \ref{thm:Stable}. 

From equation \eqref{eq:LocTrun}, we can get
$$R^n = (I - \tau \As D - \tau \Ahs L) U_*^n - \bbone u(t^n) - \tau \Ahs F^n.$$
Substitute this into the relation of $E^2$, we can get
$$E_2^n = U_*^n - (I - \tau \As D - \tau \Ahs L)^{-1} (\bbone u(t^n) + \tau \Ahs F^n).$$

The Taylor expansion gives us that
$$U_*^n = \bbone u(t^n) + \sum_{l=1}^{k} \frac{1}{l!}\tau^l c^l u^{(l)}(t^n) + O(\tau^{k+1}).$$
and 
\begin{align*}u^{(l)}(t^n) &= (D+L)u^{(l-1)}(t^n) + f^{(l-1)}(t^n) = \dots \\
    &= (D+L)^lu(t^n) + \sum_{p=0}^{l-1} (D+L)^{l-1-p}f^{(p)}(t^n). 
\end{align*}
Although $\As D + \Ahs L$ can be non-self-adjoint and unbounded, 
the laws of finite productions are true so that 
$$(I-X)^{-1} = (I-X)^{-1}(I-X^{k+1} +X^{k+1}) = (I+X+\dots +X^k) + (I-X)^{-1}X^{k+1}$$
where $X$ can be $(\tau \As D+\tau\Ahs L)$. 

Finally, Taylor expansion on $F$ is 
$$F^n = \bbone f(t^n) + \sum_{l=1}^{k-1} \frac{1}{l!}\tau^lc^lf^{(l)}(t^n) + O(\tau^k).$$

Note that we only need the last element of $E_2$, and the assumptions (P4) guarantee that
$$z^\top (\As D+\Ahs L)^l \bbone = \frac{1}{l!}(D+L)^l = \frac{1}{l!}\calA^l$$
and
$$z^\top (\As D+\Ahs L)^r\Ahs c^l = \frac{l!}{(l+1+r)!}(D+L)^r= \frac{l!}{(l+1+r)!}\calA^r.$$

Combine all the above equations, we can see that 
all the low order terms are cancelled so 
$$z^\top E_2 = C \tau^{k+1} $$
and
\begin{align*}\|e^{n+1}\| \leq& \|e^n\| + 
C\tau^{k+1} \|\calA^{k+1}u(t^n)\|_H  \\
&+C\tau^k \int_{t^n}^{t^{n+1}}\|u^{(k+1)}\|_H \rmd s 
+ C\tau^k \sum_{l=0}^{k-1} \int_{t^n}^{t^{n+1}}\|\calA^lf^{(k-l)}\|_H\rmd s
.
\end{align*}
where the constant $C$ is only related to the 
scheme itself. 

Furthermore, 
\begin{align*}
\tau\|\calA^{k+1}u(t^n)\|_H =& \int_{t^n}^{t^{n+1}}\|\calA^{k+1}u(s)\|_H\rmd s \\
&+ \int_{t^n}^{t^{n+1}}\|\calA^{k+1}u(t^n)\|_H- \|\calA^{k+1}u(s)\|_H \rmd s \\
\leq&  \int_{t^n}^{t^{n+1}}\|\calA^{k+1}u(s)\|_H\rmd s + \tau \int_{t^n}^{t^{n+1}}\|\calA^{k+1}u'(s)\|_H\rmd s \\
\end{align*}
so, finally, if $\tau \leq \tau^*$, we can get
\begin{align*}\|e^{n}\| \leq& \|e^0\| + 
    C\tau^{k} \int_{0}^{t^{n}}\|\calA^{k+1}u(s)\|_H\rmd s + C\tau^{k+1} \int_{0}^{t^{n}}\|\calA^{k+1}u'(s)\|_H\rmd s  \\
    &+C\tau^k \int_{0}^{t^{n}}\|u^{(k+1)}\|_H \rmd s 
    + C\tau^k \sum_{l=0}^{k-1} \int_{0}^{t^{n}}\|\calA^lf^{(k-l)}\|_H\rmd s
    .
    \end{align*}

If the exact solution and source term is bounded in $(0, \infty)$, then 
we can get the long time error estimate. 
Here the constant $C$ is only related to the scheme, 
and $\tau^*$ is only
related to the scheme and $\beta$ in equation \eqref{eq:beta}. 
Neither of them are related to the source term $f$,
the exact solution $u$, 
the numerical solution $u^n$, 
or the mesh size $\tau$. \qed

\begin{remark}
Note that each possible combination of assumption (P4) has appeared in the above proof during the Taylor's expansion, 
so the stage order assumption (P4) is also necessary conditions. 
\end{remark}

\section{Convergence Analysis for Semilinear IMEX-RK}

In this section, we will focusing on solving equation \eqref{eqn:AC} with IMEX-RK method. 
When we go on the nonlinear case, things are getting complicated, 
because the explicit terms, which are often the nonlinear terms, 
may be singular when they exceed their boundaries. 
One way is just like we have done before, 
that we use a extrapolation instead, 
but an extrapolation sometimes may break the stability thus 
we should pay more care on the analysis. 
The other way is to modify the potential term 
and avoid its singularity. 

To solve the equation \eqref{eqn:AC}, first, we define $\hatf$ as 
\begin{equation}\label{eq:ACmod}
    \hatf(v)= \begin{cases}
    f(v), & \text{if $\lvert v \rvert \leq \alpha$,}\\ 
    f'(\alpha)(v-\alpha), &\text{if $v > \alpha$,}\\ 
    f'(-\alpha)(v+\alpha), &\text{if $v < -\alpha$,}
\end{cases}\end{equation}

Note that $f(u) = \hatf(u)$ for $\lvert u \rvert \leq \alpha$, 
thus the equation \eqref{eqn:AC} and \eqref{eq:ACmod} share the same exact solution. 
Moreover, since the modification is tangent cutoff of the original function, 
we can know that $f \in H^2(\bbR)$. 


IMEX-RK method on equation \eqref{eqn:AC} reads as 


\begin{equation}\label{eq:RK-implicit-explicit-semilinear}
    \begin{cases}
    u^{ni} = u^{n-1} + \tau \sum_{j=0}^m a_{ij} \Delta u^{nj} + \tau \sum_{j=0}^m \hata_{ij} \hatf(u^{nj}) &\quad \text{for} ~~i=1,2,\dots,m, \\
    E(\tau) u^{n-1} = u^{n-1} + \tau \sum_{i=0}^m b_i \Delta u^{ni} + \tau \sum_{i=0}^m \hatb_i\hatf(u^{ni}), \\
    u^{n} = E(\tau) u^{n-1}
  \end{cases}
\end{equation}
which also consists an operator $u^{n} = E(\tau) u^{n-1}$. 

\begin{theorem}
    The operator $E(\tau)$ defined in equation \eqref{eq:RK-implicit-explicit-semilinear} 
    satisfies that 
    $$\|E(\tau)v - E(\tau)w\| \leq (1+C\tau) \|v-w\|$$
    for all $v, w \in L^2(\Omega)$. Here $\|\cdot\|$ refers to $L^2$ norm.  
\end{theorem}

Proof: 
Define $\bfe_i$ is the vector with $(i+1)$-th entry $1$ and others $0$. 
$U^n = \{u^{n,0}, u^{n,1}, \dots\}^\top$ , $u^{ni}= \bfe_i^\top U^n$, 
$\hatf(U^n) = \{\hatf(u^{n,0}), \hatf(u^{n,1}), \dots\}^\top$, so in vector form, 
\begin{equation}\label{eq:RK-implicit-explicit-semilinear-vector}
    U^n = \bbone u^{n-1} + \tau A \Delta U^n + \tau \hatA f(U^n),
\end{equation}
and
\begin{equation}
    U^n = (I-\tau A \Delta)^{-1}\bbone u^{n-1} + \tau (I-\tau A \Delta)^{-1} \hatA \hatf(U^n).
\end{equation}
Test with $\bfe_i$, similar to the linear case we can get
\begin{align*}
    u^{ni} & =  \bfe_i^\top(I-\tau A \Delta)^{-1}\bbone u^{n-1} + \tau  \bfe_i^\top(I-\tau A \Delta)^{-1} \hatA \hatf(U^n),\\
        & =:  \sigma_i(-\tau\Delta) u^{n-1} + \tau\sum_{j=0}^{i-1} p_{ij}(-\tau\Delta) \hatf(u^{nj}).     
\end{align*}
The boundedness of such $\sigma_i$ and $p_{ij}$ are literally the same to the linear case, 
beacuse the sub-scheme of an IMEX-RK scheme is also an IMEX-RK scheme in low order, 
and the failure of high accuracy has no influence on its stability. 

Define $v^i, w^i$ are the stages with respect to $v$ and $w$, we can get 

\begin{align*}\|v^i - w^i\| &\leq \| \sigma_i(v-w)\| + \tau\sum_{j=0}^{i-1} \left\|p_{ij} \left(\hatf(v^j)-\hatf(w^j) \right)\right\|\\
    &\leq  \|v-w\| + CL \tau\sum_{j=0}^{i-1} \left\|v^j-w^j\right\|
\end{align*}
Combine for all $i$ and dating back to the first stage, we can finally defive that 
$$\|E(\tau)(v-w)\| \leq \|v-w\|. $$
\qed

Later we will show the consistency error of semilinear IMEX-RK method. 


\begin{theorem}
    Suppose that the Assumption (P1), (P2), (P3) and (P4) are valid for $k \geq 3$, 
    $u$ is sufficiently smooth on both space and time, then  
    the operator $E(\tau)$ defined in equation \eqref{eq:RK-implicit-explicit-semilinear} 
    satisfies that 
    \begin{equation}\label{eq:semilinear-consistent}
        \|E(\tau) u(t^{n-1}) - u(t^{n})\| \leq C\tau^{k+1}
    \end{equation}
    for all $n \geq 1$. Here $\|\cdot\|$ refers to $L^2$ norm.
    Furthermore, 
    in condition that 
    $$(b_* \cdot c)^\top A_* c = \frac18$$  
    works for all $A_* \in \{A_\sigma, \hatA_\sigma\}$, 
    $b_* \in \{b_\sigma, \hatb_\sigma\}$, 
 equation \eqref{eq:semilinear-consistent} hold for $k=4$. 
 Here $ v\cdot w$ means the multiplication elementwisely. 
\end{theorem}

Proof: 

Equation \eqref{eq:RK-implicit-explicit-semilinear-vector} gives us that 
\begin{equation}\label{eq:semilinear-start}
    U^n - \bbone u(t^{n-1}) = \tau A \Delta U^n + \tau \hatA \hatf(U^n) = O(\tau).
\end{equation}
and 
\begin{equation}\label{eq:semilinear-first}
    \begin{aligned}
    U^n - \bbone u(t^{n-1}) =~& \tau \left(A  \bbone\Delta u(t^{n-1}) + \hatA \bbone \hatf(u(t^{n-1})) \right)\\
    &+ \tau A \Delta\left(U^n - \bbone u(t^{n-1})\right) + \tau \hatA \left(\hatf(U^n) - \bbone \hatf(u(t^{n-1}))\right)\\
    =~& \tau u_t^0 + O(\tau^2), 
    \end{aligned}
\end{equation}
where $u_t^0 = A \Delta \bbone u(t^{n-1}) + \hatA \bbone f(u(t^{n-1}))$. 

Substitute equation \eqref{eq:semilinear-first} into itself again, 
we can derive that 
\begin{equation}\label{eq:semilinear-second}
    \begin{aligned}
    U^n - \bbone u(t^{n-1}) =~& \tau \left(A  \bbone\Delta u(t^{n-1}) + \hatA \bbone \hatf(u(t^{n-1})) \right)\\
    &+ \tau A \Delta\left(U^n - \bbone u(t^{n-1})\right) \\
    &+ \tau \hatA \left(\hatf_u\left(U^n - \bbone u(t^{n-1})\right) + O(\tau^2)\right)\\
    =~& \tau u_t^0 + \tau A\Delta \cdot \tau u_t^0 + \tau \hatA \hatf_u \cdot \tau u_t^0  + O(\tau^3)\\
    =~& \tau u_t^0 + \tau^2 u_{tt}^0 + O(\tau^3), 
    \end{aligned}
\end{equation}
where $u_{tt}^0 = A \Delta u_t^0 + \hatA \hatf_u u_t^0$. 

Denote $v^{\cdot i}$ as the $i$-th power elementwisely and pointwisely, then
\begin{equation}\label{eq:semilinear-third}
    \begin{aligned}
    U^n - \bbone u(t^{n-1}) =~& \tau \left(A  \bbone\Delta u(t^{n-1}) + \hatA \bbone \hatf(u(t^{n-1})) \right)\\
    &+ \tau A \Delta\left(U^n - \bbone u(t^{n-1})\right) \\
    &+ \tau \hatA \left(\hatf_u\left(U^n - \bbone u(t^{n-1})\right) 
     + \frac12 \hatf_{uu}\left(U^n - \bbone u(t^{n-1})\right)^{\cdot2}  + O(\tau^3)\right)\\
    =~& \tau u_t^0 + \tau A\Delta \cdot \left(\tau u_t^0 + \tau^2 u_{tt}^0\right) \\
     & + \tau \hatA \hatf_u \cdot\left(\tau u_t^0 + \tau^2 u_{tt}^0\right)
     + \tau \hatA \frac12 \hatf_{uu} (\tau u_t^0)^{\cdot 2}  + O(\tau^4)\\
    =~& \tau u_t^0 + \tau^2 u_{tt}^0+ \tau^3 u_{ttt}^0 + O(\tau^4), 
    \end{aligned}
\end{equation}
where $u_{ttt}^0 = A \Delta u_{tt}^0 + \hatA \hatf_u u_{tt}^0 + \frac12 \hatA \hatf_{uu} (u_t^0)^{\cdot 2}$. 
Similarily, we reached that 
\begin{equation}\label{eq:semilinear-fourth}
    \begin{aligned}
    U^n - \bbone u(t^{n-1}) =  \tau u_t^0 + \tau^2 u_{tt}^0+ \tau^3 u_{ttt}^0+ \tau^4 u_{tttt}^0 + O(\tau^5), 
    \end{aligned}
\end{equation}
where $u_{tttt}^0 = A \Delta u_{ttt}^0 + \hatA \hatf_u u_{ttt}^0 + \hatA \hatf_{uu} u_t^0 \cdot u_{tt}^0 + \frac16 \hatA (u_t^0)^{\cdot 3}$. 
Test the above equations with $\bfe_m$, and we can get
$$E(\tau) u(t^{n-1}) = u(t^{n-1}) + \bfe_m^\top \left(\tau u_t^0 + \tau^2 u_{tt}^0+ \tau^3 u_{ttt}^0+ \tau^4 u_{tttt}^0\right) + O(\tau^5),$$
where the RHS is only depend on $u(t^{n-1})$. 
Assume that the exact solution $u$ is smooth enough, and compare it with the Taylor's expansion, 
we can find that all the low-order terms disappeared, thus the theorem is proved. \qed

After we have proved the consistency, the final convergence theorem is straightforward. 

\begin{theorem}
    Suppose that the Assumption (P1), (P2), (P3) and (P4) are valid for $k \geq 4$, 
    and $(b_* \cdot c)^\top A_* c = \frac18 $  
    works for all $A_* \in \{A_\sigma, \hatA_\sigma\}$, 
    $b_* \in \{b_\sigma, \hatb_\sigma\}$ if $k=4$.  
    The exace solution $u$ of equation \eqref{eqn:AC} is sufficiently smooth 
    on both space and time $u^n$ is the solution of \eqref{eq:RK-implicit-explicit-semilinear}, 
    then
    \begin{equation}\label{eq:semilinear-convergence}
        \|u^{n} - u(t^{n})\| \leq C\tau^{k}.
    \end{equation}

\end{theorem}
Prove: 
\begin{equation}\label{eq:semilinear-convergence-proof}
    \begin{aligned}
    \|u^{n} - u(t^{n})\| &= \|E(\tau)u^{n-1} - u(t^{n})\|\\
    &\leq  \|E(\tau)u^{n-1} - E(\tau)u(t^{n-1})\| + \|E(\tau)u(t^{n-1}) - u(t^{n})\|\\
    &\leq  (1+C\tau) \|u^{n-1}- u(t^{n-1})\| + C\tau^{k+1} .
    \end{aligned}
\end{equation}
Then Gr\"onwall's inequality will finish the proof. 

\begin{remark}
    Since this is a single-step method, as what we have done in equation \eqref{eqn:semi-cut}, 
    we can build a scheme that
    \begin{equation} \label{eq:semilinear-cutoff}
        u^{n} =  \min(\max( E(\tau) u^{n-1},-\alpha),\alpha),
    \end{equation}

    The exact solution will always meet the maximum bound condition, 
    so $$\|u^{n} - u(t^{n})\| \leq \|E(\tau)u^{n-1} - u(t^{n})\|.$$
    The rest part of the proof is the same to \eqref{eq:semilinear-convergence-proof}. 
\end{remark}

\begin{remark}\label{Remark:IMEX-Energy}
    This analysis is compatible with the work of the paper by Fu and Yang, 
    in which they proof that 
    some IMEX-RK scheme can keep energy dissipation law with a stabilizer, 
    and we have found the scheme of $1$, $2$ and $3$ order, 
    which means that we have found a third order IMEX-RK scheme that can keep 
    both maximum bound principle and energy dissipation law.  
\end{remark}

